\subsection{Graph Neural Networks}
\label{sec:prelim:gnn}

\subsubsection{Message Passing Framework}
\label{sec:prelim:gnn:mp}
A Graph Neural Network (GNN) computes a representation $h_v^{l} \in \mathbb{R}^{w}$ for
every node $v$ by repeated neighbourhood aggregation
\citep{kipf2017gcn,hamilton2017graphsage}. Writing $\mathcal{N}(v)$ for the neighbours of
$v$ and $e_{uv}$ for the attribute of edge $(u, v)$, one layer computes
\begin{equation}
    h_v^{l} = \phi^{l}\!\left( h_v^{l-1},
        \bigoplus_{u \in \mathcal{N}(v)} \psi^{l}\!\left( h_u^{l-1}, e_{uv} \right)
    \right),
    \label{eq:prelim:mp}
\end{equation}
where $\psi^{l}$ builds a message per edge, $\bigoplus$ is a permutation-invariant
aggregation such as a sum, and $\phi^{l}$ updates the node state. The encoder of this
thesis is one instantiation of~\eqref{eq:prelim:mp}, written out
in~\ref{sec:method:architecture:equations}.

After $L$ layers, $h_v^{L}$ is a function of exactly the $L$-hop neighbourhood of $v$ and
of nothing outside it. On an AIG that neighbourhood is a depth-$L$ prefix of the cones
of~\eqref{eq:prelim:cone}, so an encoder of depth $L$ reads $L$ levels of logic around each
gate. Every depth-coupling argument in this thesis rests on this one fact: the refinement
depth of~\ref{sec:method:summary:wl} is set to $L$, the receptive-field metric
of~\ref{sec:method:experiment:metrics:reduction} counts the gates inside $L$ hops, and
hypothesis H1 (\ref{sec:intro:rq3}) is a claim about what path contraction does to those
$L$ hops.

\subsubsection{Graph-Level Readout and Pooling}
\label{sec:prelim:gnn:readout}
A graph-level prediction pools the node representations into one vector, here by a mean
over nodes, optionally after summing the per-layer outputs so that shallow and deep
features both reach the readout (jumping knowledge, \citealp{xu2018jk}). Pooling imposes
the constraint that shapes the whole experimental setup: every node of a graph must be
present in the same forward pass, so a graph is the indivisible unit of batching and
cannot be split across steps. Batches are therefore assembled under a node budget rather
than as a fixed graph count (\ref{sec:method:experiment:batching}), and the memory cost
of one step is set by the largest graphs, which is exactly the cost reduction is meant to
lower.

\subsubsection{Positional and Structural Encodings}
\label{sec:prelim:gnn:pe}
Message passing sees only relative structure: two nodes whose $L$-hop neighbourhoods are
isomorphic receive identical representations wherever they sit in the graph. A Positional
Encoding (PE) breaks that indifference by attaching a coordinate to each node. On an
undirected graph no canonical coordinate exists and encodings are built from random walks
or spectra; on a DAG the topological level $\ell(v)$ of~\eqref{eq:prelim:level} is an
absolute coordinate that the graph carries for free. The level-based encoding used here
(\ref{sec:method:data:representation:pe}) is therefore a DAG-native choice rather than a
generic one.

\subsubsection{Expressivity and the Weisfeiler-Leman Hierarchy}
\label{sec:prelim:gnn:expressivity}
Colour refinement, equivalently the one-dimensional Weisfeiler-Leman algorithm (1-WL),
iteratively refines a node colouring by the multiset of neighbour colours,
\begin{equation}
    c_v^{(t)} = \operatorname{hash}\!\left( c_v^{(t-1)},\
        \bigl\{\!\!\bigl\{\, c_u^{(t-1)} : u \in \mathcal{N}(v) \,\bigr\}\!\!\bigr\}
    \right),
    \label{eq:prelim:cr}
\end{equation}
starting from initial labels $c_v^{(0)}$. A partition of $\mathcal{V}$ is
\emph{equitable} when every node of a class has the same number of neighbours in every
class; run to a fixed point,~\eqref{eq:prelim:cr} computes the coarsest equitable
partition refining the initial colouring \citep{grohe2014colourrefinement}. The
connection to~\eqref{eq:prelim:mp} is an upper bound: a message-passing network is at
most as discriminative as 1-WL, so two nodes sharing a colour after $L$ rounds receive
identical representations from every $L$-layer encoder \citep{xu2019gin}.

Replacing the multiset in~\eqref{eq:prelim:cr} by a set yields bisimulation, which ignores
how many neighbours of a colour exist and asks only whether one does. A count cap $c$
interpolates between the two, truncating each multiplicity at $c$: the cap $c = 1$ is
bisimulation and $c = \infty$ is colour refinement \citep{bollen2023exact}. The bound
above is what makes quotienting by such a partition lossless rather than approximate:
merging nodes the network provably cannot distinguish removes no information the network
could have used, and the automorphism orbits of a graph form an equitable partition, so
genuinely symmetric structure is always compressible this way. AIGs have such structure
in quantity. The two fanins of an AND gate are interchangeable, datapath logic replicates
one bit slice across the word width, and standard sub-circuits recur throughout a design,
each of which creates orbits that a random graph does not have. How much of this
redundancy survives structural hashing is measured in~\ref{sec:method:summary:wl}.

\subsubsection{Training Cost, Memory, and Inference}
\label{sec:prelim:gnn:cost}
Training and inference differ in what must be held in memory. Backpropagation stores the
activations of every layer for every node in the batch, so training memory grows with
$L \cdot |\mathcal{V}|$ on top of parameters and optimizer state, and the largest graph in
the corpus sets the peak. Forward-only inference stores none of this. The asymmetry is
the mechanism RQ5 (\ref{sec:intro:rq5}) depends on: a model that must be trained on
reduced graphs to fit in device memory may still be queried on full graphs afterwards.

Four techniques recur in Chapter 3 and are defined here once. Reported memory is
\emph{allocated} memory, the tensors actually held, as opposed to the larger
\emph{reserved} pool the allocator keeps (\ref{sec:method:experiment:metrics:efficiency}).
Mixed precision stores activations in half-width floating point and roughly halves their
cost. Gradient checkpointing discards intermediate activations and recomputes them during
the backward pass, trading compute for memory, and is what makes the largest baseline
runnable at all (\ref{sec:method:architecture:baselines}). Gradient accumulation sums
gradients over several small steps before updating, so the effective batch size exceeds
what fits in one step. All four trade something for memory at a \emph{fixed} input size,
which is what separates them from the input reduction this thesis studies
(\ref{sec:relwork:scaling:memory}).

Depth has a cost of its own. A fixed-width representation cannot carry the information of
a receptive field that grows exponentially with depth, so signal from distant nodes is
compressed away, a failure known as over-squashing \citep{alon2021oversquashing}.
Contracting paths shortens the distance signal must travel and is a recognised remedy
\citep{dwivedi2022lrgb}. That contraction is the mechanism behind hypothesis H1
(\ref{sec:intro:rq3}), which expects summarization to help where sparsification can only
remove.
